\section{Design and Debugging}

\subsection{Defects, Errors, and Failures}

\begin{definition}
    \textbf{Defect, Error, and Failure Taxonomy:}
    \begin{itemize}
        \item \textbf{Defect (Fault / Bug):} An imperfection, flaw, or mistake in the source code or design created by the programmer (e.g., an off-by-one loop condition, an uninitialized variable, or a missing boundary check).
        \item \textbf{Error:} An erroneous or invalid internal runtime state of the program resulting from the execution of a defect (e.g., an internal variable holding an illegal value or corrupted memory on the heap). An error is internal and may not yet be observable from outside the system.
        \item \textbf{Failure:} An externally observable divergence of the system's behavior from its specification or expected behavior (e.g., a segmentation fault / crash, incorrect output displayed to the user, an infinite loop, or data corruption).
    \end{itemize}
\end{definition}

\begin{remark}
    \textbf{The Fault Propagation Chain:}
    $$\text{Defect (Code)} \xrightarrow{\text{Execution}} \text{Error (State)} \xrightarrow{\text{Propagation}} \text{Failure (Symptom)}$$
\end{remark}

\subsection{Testing Principles}

\begin{definition}
    \textbf{Core Testing Strategies:}
    \begin{itemize}
        \item \textbf{Test Early:} Testing should begin alongside development rather than being deferred until the end. Finding and resolving defects early reduces debugging complexity and cost.
        \item \textbf{Test Incrementally:} Build and test software in small, verifiable steps. When code is tested after every small addition, newly introduced bugs are localized to the most recent changes.
        \item \textbf{Test Small, Simple Units:} Isolate individual functions and components (unit testing) to verify their behavior independently before combining them into complex systems.
        \item \textbf{Test Boundary Conditions (Corner Cases):} Software bugs disproportionately cluster around edges and extreme limits. Thorough test suites must evaluate:
    \end{itemize}
\end{definition}

\begin{remark}
    \textbf{Dijkstra's Aphorism on Testing:}
    \begin{quote}
        \emph{``Program testing can be used to show the presence of bugs, but never to show their absence!''}
        \hfill --- Edsger W. Dijkstra (1970)
    \end{quote}
    No matter how many test cases pass, testing cannot formally prove that a program is defect-free because testing all possible inputs and system states is computationally infeasible. Testing must be combined with disciplined design, modularity, assertions, and invariant verification.
\end{remark}

\subsection{Debugging Advice and Strategies}

\begin{definition}
    \textbf{Systematic Debugging Guidelines:}
    \begin{itemize}
        \item \textbf{Examine the Most Recent Change:} If a previously working system breaks, the bug almost certainly lies within the most recent modifications. Use version control tools (\texttt{git diff}, \texttt{git log -p}, \texttt{git bisect}) to isolate changes.
        \item \textbf{Simplify the Input (Minimal Reproducible Example):} Reduce large, complex failing inputs to the smallest, simplest possible test case that still triggers the failure. Eliminating irrelevant operations simplifies tracing program state.
        \item \textbf{Inspect the Stack Trace:} Upon a crash or assertion failure, use a debugger (e.g., \texttt{gdb} command \texttt{backtrace} or \texttt{bt}) or core dump to identify the exact call chain, active stack frames, and parameter values at the point of failure.
        \item \textbf{Sanity Check Intermediate Results:} Formulate concrete hypotheses about program state and verify them at key execution milestones using assertions (\texttt{assert()}), logging, or debugger watchpoints. Verify preconditions, postconditions, and loop invariants.
        \item \textbf{Understand the Root Cause:} Avoid superficial "band-aid" fixes (such as blindly changing \texttt{<} to \texttt{<=}, adding arbitrary \texttt{+1} offsets, or silencing warnings). Only modify code once the exact failure mechanism and the rationale for the fix are fully understood.
    \end{itemize}
\end{definition}

\begin{remark}
    \textbf{Essential Systems Debugging Tools:}
    \begin{itemize}
        \item \textbf{GDB (GNU Debugger):} Interactive debugger to step through execution (\texttt{step}, \texttt{next}), inspect variable values (\texttt{print}), set breakpoints/watchpoints, and inspect crash backtraces (\texttt{bt}).
        \item \textbf{Valgrind (Memcheck):} Dynamic memory analysis tool detecting memory leaks, invalid heap reads/writes, and use-after-free bugs.
        \item \textbf{Sanitizers (GCC / Clang):} Compile-time instrumentation adding automated checks:
        \begin{itemize}
            \item AddressSanitizer (\texttt{-fsanitize=address}): Catches buffer overflows and use-after-free bugs.
            \item UndefinedBehaviorSanitizer (\texttt{-fsanitize=undefined}): Catches undefined behavior (e.g., signed overflow, invalid shifts, null dereferences).
        \end{itemize}
    \end{itemize}
\end{remark}
